% Nota histórica: de dónde sale la épsilon (castellano).

\begin{frame}
\didactatitle{Una breve historia de la épsilon}

Desde Newton y Leibniz, a finales del siglo \textsc{xvii}, se calcularon
límites durante más de un siglo sin definirlos: se razonaba con cantidades
«evanescentes» o «infinitamente pequeñas», y funcionaba, pero nadie sabía
decir con precisión por qué.

\begin{itemize}
  \item \textbf{1817.} Bolzano demuestra el teorema de los valores intermedios
        sin apoyarse en la intuición geométrica.
  \item \textbf{1821.} Cauchy, en su \emph{Cours d'analyse}, define el límite
        y la continuidad con palabras: una variable que se acerca
        indefinidamente a un valor fijo.
  \item \textbf{Hacia 1860.} Weierstrass, en sus clases de Berlín, da a la
        definición la forma con $\varepsilon$ y $\delta$ que usamos hoy.
\end{itemize}

\onlynotes{%
  Lo que cambió no fue el cálculo, sino la manera de estar seguros de él: con
  $\varepsilon$ y $\delta$, «se acerca» deja de ser una imagen y pasa a ser
  una desigualdad.
}
\end{frame}
